\begin{lemma}[Ordering of supporting and extension vertices]
Let \(T\) be a positive-rung extended Tashkinov tree generated by
\[
 \mathcal T=\{(T_i,\varphi_{i-1},S_{i-1},F_{i-1},\Theta_{i-1}):
                   1\leq i\leq n+1\},
\]
and suppose that \(T\) has the maximum property. Keep the graph \(G\),
palette size \(k\), and distinguished uncolored edge \(e_0\) fixed, but
allow the certified input coloring to vary arbitrarily. Assume explicitly
the following six assertions for every positive-rung ETT \(T'\) on these
fixed data, with arbitrary certified input coloring \(\varphi'_0\), maximum
property, and rung number \(s<n\), writing \(\varphi'_s\) for its generating
coloring and \(\mathcal T'\) for its certified series:
\begin{enumerate}
\item \(V(T')\) is elementary with respect to \(\varphi'_s\).
\item Under every proper \((T'_{s+1},D'_s,\varphi'_s)\)-stable coloring
leaving precisely \(e_0\) uncolored, \(T'_{s+1}\) is interchangeable for
any pair of colors of which at least one is missing on \(T'_{s+1}\).
\item For each PE slot \(a\leq s\) with supporting vertex \(u_a\) and
\(m(u_a)=b\), every proper \((T'_a,D'_a,\varphi'_a)\)-stable coloring
leaving precisely \(e_0\) uncolored is both
\((T'(u_a)-u_a,D'_{b-1},\varphi'_{b-1})\)-stable and
\((T'_{b-1},D'_{b-1},\varphi'_{b-1})\)-stable. Moreover, the stored
color pairs at two distinct PE slots with distinct supporting vertices
and the same first-appearance index are disjoint.
\item If the terminal slot is PE with supporting vertex \(u_s\) and
stored colors \(\gamma'_s,\delta'_s\), then under every proper
\((T'_s,D'_s,\varphi'_s)\)-stable coloring leaving precisely \(e_0\)
uncolored, the \((\gamma'_s,\delta'_s)\)-component through \(u_s\)
contains exactly the one vertex \(u_s\) of \(T'_s\).
\item If \(\tau\) is any proper \((T'_s,D'_s,\varphi'_s)\)-stable
coloring leaving precisely \(e_0\) uncolored, \(d\) is defective on
\(T'_s\) under \(\tau\), and \(x\) is a nonminimum vertex of
\(I[\partial_{\tau,d}(T'_s)]\), then
\(x\preceq_{T'}u_a\) for every PE supporting vertex or SE extension
vertex \(u_a\) with \(m(x)\leq a\leq s\).
\item Every proper \((T'_s,D'_s,\varphi'_s)\)-stable coloring leaving
precisely \(e_0\) uncolored is a \(\varphi'_s\bmod T'_s\) coloring.
\end{enumerate}

Write \(T_n\) for the terminal pre-extension tree of \(\mathcal T\).
Let \(\sigma_n\) be any proper
\((T_n,D_n,\varphi_n)\)-stable coloring leaving precisely \(e_0\)
uncolored, let \(d\) be defective on \(T_n\) under \(\sigma_n\), and let
\(v\) be a nonminimum vertex of
\(I[\partial_{\sigma_n,d}(T_n)]\). If \(v_i\) is the exact chosen
vertex at a PE or SE slot \(i\leq n\), and \(m(v)\leq i\), then
\[
                              v\preceq_T v_i .
\]
Here \(v\preceq_Tv_i\) means \(v=v_i\) or \(v\prec_Tv_i\) in the exact
vertex order of the displayed final tree \(T\). No chosen vertex at an
RE slot is asserted to be a supporting or extension vertex.
\end{lemma}
\begin{proof}
We first discharge the two current-rung assertions which the manuscript
labels (4.3), rather than treating either of them as an extra hypothesis.
Specialize the six lower-rung assumptions to certified input
\(\varphi'_0=\varphi_0\). They then meet the lower-rung premise of
\noderef{ExtensionRules}; hence current assertion (iii), including both
stability conclusions at every current PE slot and the disjointness
conclusion, holds for \(\mathcal T\). The original generalized package,
with its arbitrary input-coloring quantifier still present, meets the
premise of \noderef{ExtensionBase}. Thus current assertion (iv) also
holds: at a terminal PE slot the component of the stored gamma and delta
through its supporting vertex meets \(V(T_n)\) only at that vertex.

Fix \(\sigma_n,d,v,i,v_i\) as in the statement. Put \(X=V(T_n)\).
Choose an exact \(d\)-colored boundary identity \(a\) incident with \(v\),
and choose an earlier vertex \(w\prec_{T_n}v\) and an exact
\(d\)-colored boundary identity \(b\) incident with \(w\). These
witnesses exist by \noderef{BoundaryColorVertex}. They are distinct.
Indeed, \noderef{BoundaryEdges} gives an inside and an outside endpoint
for a boundary identity, whereas \(w,v\in X\) and \(w\ne v\) by
\noderef{TreeVertexOrder}; a single boundary identity cannot therefore
be incident with both \(w\) and \(v\). The membership and color equations
for \(a,b\) are exact edge-identity statements by
\noderef{BoundaryColorEdges}. We retain these identities throughout.

The certificate \noderef{TashkinovSeriesStep} at the terminal positive
slot has exactly one of the PE, SE, and RE constructors. We treat PE
first and then SE/RE.

\textbf{Case 1: the terminal slot is PE.}
Let its exact maximum-defective choice be
\((\pi_{n-1},v_n,f_n,u_n,\delta_n)\), let its second stable coloring be
\(\pi'_{n-1}\), and let its PE color be \(\gamma_n\). The PE fields in
\noderef{TashkinovSeriesStep} identify the chosen vertex and stored fields
with \(v_n,\gamma_n,\delta_n\), identify \(F_n=\{f_n\}\), and state that
\(\varphi_n\) is the switch of the full
\((\gamma_n,\delta_n)\)-component of \(v_n\) under
\(\pi'_{n-1}\). The choice data in
\noderef{MaximumDefectiveChoice} say that \(\pi_{n-1}\) is proper,
leaves precisely \(e_0\) uncolored, is
\((T_n,D_{n-1},\varphi_{n-1})\)-stable, and colors the exact incident
identity \(f_n\) by \(\delta_n\). The PE constructor says that
\(\pi'_{n-1}\) is proper, leaves precisely \(e_0\) uncolored, and is
stable relative to \(\pi_{n-1}\) with control
\(D_{n-1}\cup\{\delta_n\}\).

We record two exact consequences at \(v_n\). The rule
\noderef{ParallelExtensionColor} makes \(\gamma_n\) missing at \(v_n\)
under \(\pi_{n-1}\), and the vertex clause of
\noderef{StableColoring} transfers this to \(\pi'_{n-1}\). Since
\(\delta_n\) belongs to the enlarged control set, the edge clause fixes
\(\pi'_{n-1}(f_n)=\delta_n\). By \noderef{AdjacentEdges}, every other
identity incident with \(v_n\) is adjacent to \(f_n\), so properness in
\noderef{IsProperEdgeColoring} makes \(f_n\) the unique
\(\delta_n\)-colored identity incident with \(v_n\), and missingness says
there is no incident \(\gamma_n\)-identity. Hence the switch in
\noderef{TwoColorComponentSwap} changes \(f_n\) from \(\delta_n\) to
\(\gamma_n\), changes no other incident identity to \(\delta_n\), and
therefore makes
\[
                 \delta_n\in\overline\varphi_n(v_n).           \tag{1}
\]
The vertex clause of the assumed terminal stability now gives
\[
                 \delta_n\in\overline\sigma_n(v_n).           \tag{2}
\]

Next we identify the gamma boundary edge, in both directions and by
identity. The first conclusion of
\noderef{HkuParallelBoundaryStructure}, applied to the exact terminal PE
fields, is
\[
                \partial_{\varphi_n,\gamma_n}(X)=\{f_n\}.      \tag{3}
\]
The terminal stored-set equation includes \(\gamma_n\) in \(S_n\).
By the set-difference definition \noderef{TashkinovStoredColors}, either
\(\gamma_n\) is missing on \(T_n\), or it belongs to \(D_n\); in either
case it belongs to
\(\overline\varphi_n(T_n)\cup D_n\). Thus the edge clause of terminal
stability fixes every identity in the left side of (3). Conversely, if
an identity has \(\sigma_n\)-color \(\gamma_n\), the symmetric stability
criterion \noderef{sc1}, using properness and exact
one-edge-uncoloredness of both colorings, fixes it back to its
\(\varphi_n\)-value. Boundary membership depends only on its exact
endpoints. Consequently
\[
                \partial_{\sigma_n,\gamma_n}(X)=\{f_n\}.       \tag{4}
\]
Since \(d\) is defective, \noderef{DefectiveColor} says its boundary has
at least two identities. Equation (4) therefore gives
\[
                              d\ne\gamma_n.                     \tag{5}
\]

Apply the current assertion (iv) obtained from
\noderef{ExtensionBase}. For the component
\(K=P_{v_n}(\gamma_n,\delta_n,\sigma_n)\), it gives the exact
intersection
\[
                x\in X\ \hbox{and}\ x\in K
                       \quad\Longleftrightarrow\quad x=v_n.    \tag{6}
\]
All hypotheses of \noderef{ParallelExtensionUndoStable} are now met.
Define the recovered coloring by the edge-identity switch
\[
             \tau=\sigma_n/P_{v_n}(\gamma_n,\delta_n,\sigma_n).
\]
That node gives, in particular, that \(\tau\) is proper, leaves precisely
\(e_0\) uncolored, and is
\((T_n,D_{n-1},\varphi_{n-1})\)-stable.                         \tag{7}

We prove the boundary transport needed below, instead of merely asserting
it. Let \(g\) be any exact identity in
\(\partial_{\sigma_n,d}(X)\), and let \(x\in X\) be its inside endpoint,
as supplied by \noderef{BoundaryEdges}. Suppose that the edge-node \(g\)
belonged to \(K\). The incidence rule in
\noderef{TwoColorComponent} then also puts the adjacent vertex-node \(x\)
in \(K\), so (6) gives \(x=v_n\). Because \(g\) has
\(\sigma_n\)-color \(d\) and is an edge-node of the two-color component,
its color is either \(\gamma_n\) or \(\delta_n\). In the first
alternative (4) identifies \(g\) with the unique gamma identity \(f_n\)
and gives \(d=\gamma_n\), contrary to (5). In the second alternative
\(g\) is a \(\delta_n\)-colored identity incident with \(v_n\), contrary
to the missingness (2), expanded using \noderef{MissingColors}. Thus
no \(d\)-colored boundary identity lies in \(K\).

The switch definition \noderef{TwoColorComponentSwap} consequently
leaves every such identity \(g\) colored \(d\). Its endpoints do not
change, so it remains in the same boundary. We have proved the
identity-wise inclusion
\[
          \partial_{\sigma_n,d}(X)
             \subseteq\partial_{\tau,d}(X).                   \tag{8}
\]
In particular, the distinct identities \(a,b\) and their inside endpoints
\(v,w\) survive. Thus \(d\) is defective under \(\tau\), \(v\) is a
member of \(I[\partial_{\tau,d}(X)]\), and the same exact witness
\(w\prec_{T_n}v\) makes it nonminimum. Expanding
\noderef{DefectiveEdge} and \noderef{DefectiveVertex}, \(v\) is therefore
a defective vertex of \(T_n\) under \(\tau\).

Suppose first that the requested slot is terminal, \(i=n\). Since the
exact terminal kind is PE, the applicability disjunction forces its PE
alternative, and injectivity of the chosen-vertex equation gives
\(v_i=v_n\). Apply the maximality field of
\noderef{MaximumDefectiveChoice} to the admissible coloring \(\tau\) from
(7) and its defective vertex \(v\). It gives
\(v=v_n\) or \(v\prec_{T_n}v_n\). The terminal tree is an ordered initial
segment of the final ETT by \noderef{ExtendedTashkinovTree}; the
entrywise equalities in \noderef{TreeSequenceInitialSegment} transport
the strict alternative to \(v\prec_Tv_n\). Hence
\(v\preceq_Tv_i\), as required.

It remains in this PE case to assume \(i<n\). The inequality
\(m(v)\leq i\), together with the least-stage definition
\noderef{ETTFirstAppearanceIndex} and the nested ordered stages in
\noderef{TashkinovSeriesStep}, puts \(v\) in \(V(T_{n-1})\). Since
\(T_{n-1}\) is an ordered initial segment of \(T_n\) and contains the
displayed vertex \(v\), vertex injectivity in \noderef{TreeSequence} and
the entrywise prefix in \noderef{TreeSequenceInitialSegment} show that
the prefix contains every vertex before
\(v\), in particular \(w\); it also retains
\(w\prec_{T_{n-1}}v\). Let the outside endpoints of \(a,b\) relative to
\(X\) be \(x_a,x_b\). They lie outside \(X\), hence outside the smaller
\(V(T_{n-1})\). The same exact endpoint equations therefore put the
distinct identities \(a,b\) in
\(\partial_{\tau,d}(V(T_{n-1}))\), by
\noderef{BoundaryEdges} and \noderef{BoundaryColorEdges}. Thus \(d\) is
defective on \(T_{n-1}\), and \(v\) is still a nonminimum vertex of its
\(d\)-boundary incidence set.

Restrict \(\mathcal T\) through the preceding iteration. The certified
steps and strict ordered extension in \noderef{TashkinovSeriesStep},
together with \noderef{ExtendedTashkinovTree}, give a positive-rung ETT
\(B\) with \(n-1\) rungs, final ordered tree \(T_n\), terminal tree
\(T_{n-1}\), generating coloring \(\varphi_{n-1}\), and original input
\(\varphi_0\). Positivity follows from \(1\leq i<n\). The initial
maximum clause and all stage clauses below \(n-1\) restrict directly from
\noderef{MaximumProperty}, so \(B\) has maximum property. Relation (7)
restricts to
\((T_{n-1},D_{n-1},\varphi_{n-1})\)-stability: its vertex equalities
restrict to the smaller vertex set, and any color missing on the smaller
tree has the same vertex witness on \(T_n\), so its controlled-edge clause
also restricts by \noderef{StableColoring}.

Apply lower-rung assertion (v) from the local six-assertion hypothesis to
this exact ETT \(B\), coloring \(\tau\), color \(d\), vertices \(w,v\),
and the unchanged applicable slot \(i\). The finite minimum in
\noderef{ETTFirstAppearanceIndex} is unchanged for \(v\): its least stage
is at most \(i<n\), and the restricted and original series agree at every
candidate through that stage. Hence its required first-appearance
inequality is precisely \(m(v)\leq i\). The lower assertion yields
\(v=v_i\) or \(v\prec_{T_n}v_i\). The same initial-segment transport from
\(T_n\) to \(T\) yields \(v\preceq_Tv_i\). This completes the PE case.

\textbf{Case 2: the terminal slot is SE or RE.}
We first derive the old-reference stability used in both subcases. The
first inclusion of \noderef{HkuMissingStoredMonotone} is
\[
 \overline\varphi_{n-1}(T_n)\cup D_{n-1}
       \subseteq \overline\varphi_n(T_n)\cup D_n.              \tag{9}
\]
The terminal vertex clause of \noderef{StableColoring} identifies the
missing sets of \(\sigma_n\) and \(\varphi_n\) on \(T_n\), while its edge
clause and (9) allow the control set to be reduced from \(D_n\) to
\(D_{n-1}\). Therefore \(\sigma_n\) is
\((T_n,D_{n-1},\varphi_n)\)-stable. In the SE constructor,
\noderef{TashkinovSeriesStep} identifies \(\varphi_n\) with the selected
coloring in \noderef{MaximumDefectiveChoice}, whose selected-stability
field makes it \((T_n,D_{n-1},\varphi_{n-1})\)-stable. In the RE
constructor the step gives the exact equality
\(\varphi_n=\varphi_{n-1}\), so the same stability is reflexive. In
either case, transitivity from \noderef{sc2} gives
\[
       \sigma_n\text{ is }
       (T_n,D_{n-1},\varphi_{n-1})\text{-stable}.               \tag{10}
\]

If \(i=n\), the applicability disjunction rules out RE and forces the SE
constructor. Its chosen-vertex equation gives \(v_i\) as the selected
maximum-defective vertex. The defective color \(d\) and the exact
boundary identity \(a\) make \(v\) a defective vertex by
\noderef{DefectiveEdge} and \noderef{DefectiveVertex}. Apply the
maximality field of \noderef{MaximumDefectiveChoice} to \(\sigma_n\),
using (10): it gives \(v=v_i\) or \(v\prec_{T_n}v_i\). Ordered-prefix
transport to \(T\), exactly as in Case 1, proves the required conclusion.
For a terminal RE the same applicability equation is contradictory, so
there is no terminal RE subcase.

Finally suppose \(i<n\). The first-appearance and ordered-prefix argument
from Case 1 again puts both \(v\) and its earlier witness \(w\) in
\(T_{n-1}\). The outside endpoints of the exact identities \(a,b\)
remain outside the smaller tree, so those two distinct identities prove
that \(d\) is defective on \(T_{n-1}\) and that \(v\) is a nonminimum
member of its \(d\)-boundary incidence set. Stability (10) restricts to
\((T_{n-1},D_{n-1},\varphi_{n-1})\)-stability by the same vertex-witness
argument. Use the same positive \((n-1)\)-rung restriction \(B\), whose
maximum property and first-appearance agreement were verified above, and
apply lower-rung assertion (v). It gives
\(v=v_i\) or \(v\prec_{T_n}v_i\), and the exact ordered inclusion of
\(T_n\) in \(T\) gives \(v\preceq_Tv_i\). This proves every PE or SE
applicable slot and completes the assertion.
\end{proof}
